% ============================================================
% DOCUMENTAÇÃO DO PRIMEIRO MVP — PAC VI — TORQUE GESTÃO
% Católica de Santa Catarina
% ============================================================
% chktex-file 1
% chktex-file 13
% chktex-file 44

\documentclass[11pt, a4paper]{article}

\usepackage[T1]{fontenc}
\usepackage[utf8]{inputenc}
\usepackage[brazilian]{babel}
\usepackage{times}
\usepackage[top=2.5cm, bottom=2.5cm, left=2.5cm, right=2.5cm]{geometry}
\usepackage{setspace}
\usepackage{indentfirst}
\usepackage{titlesec}
\usepackage{xcolor}
\usepackage{fancyhdr}
\usepackage{enumitem}
\usepackage{array}
\usepackage{longtable}
\usepackage{float}
\usepackage{ragged2e}
\usepackage[hidelinks]{hyperref}

\definecolor{torquenavy}{HTML}{1B2B4B}
\definecolor{torqueamber}{HTML}{F0A500}

\setstretch{1.25}
\setlength{\parindent}{1.1cm}
\setlength{\parskip}{3pt}
\setlength{\emergencystretch}{2em}

\newcolumntype{L}[1]{>{\RaggedRight\arraybackslash}p{#1}}

\pagestyle{fancy}
\fancyhf{}
\fancyhead[R]{\thepage}
\renewcommand{\headrulewidth}{0pt}

\titleformat{\section}{\normalfont\bfseries\large\color{torquenavy}}{\thesection}{1em}{}
\titleformat{\subsection}{\normalfont\bfseries\color{torquenavy}}{\thesubsection}{1em}{}
\titleformat{\subsubsection}{\normalfont\bfseries\itshape}{\thesubsubsection}{1em}{}

\begin{document}

% --- CABEÇALHO INSTITUCIONAL ---
\thispagestyle{empty}
\begin{center}
    {\small CENTRO UNIVERSITÁRIO --- CATÓLICA DE SANTA CATARINA}\\
    {\small ENGENHARIA DE SOFTWARE --- PAC VI}\\
    {\small Documentação do Primeiro MVP do PAC VI}\\[14pt]
    {\LARGE\bfseries Documentação do Primeiro MVP}\\[6pt]
    {\Large\itshape Torque Gestão --- Sistema de Gestão Automotiva}
\end{center}

\vspace{8pt}
{\color{torqueamber}\rule{\linewidth}{2.5pt}}
\vspace{18pt}

% --- FICHA DE IDENTIFICAÇÃO ---
\renewcommand{\arraystretch}{1.35}
\noindent
\begin{tabular}{|L{3.2cm}|L{12.3cm}|}
\hline
\textbf{Nome do projeto} & Torque Gestão --- Sistema de Gestão Automotiva \\
\hline
\textbf{Integrantes} & Miguel Angelo Dufloth Filho; Miguel Angel Balladares Huertas; Leonardo Lotério de Lima; Lucas Honorato dos Santos \\
\hline
\textbf{Repositório} & \url{https://github.com/miguettohuertas/torque-gestao} \\
\hline
\textbf{Protótipo navegável} & \url{https://torque-gestao.surge.sh} \\
\hline
\textbf{Data de elaboração} & Joinville, 17 de agosto de 2026 \\
\hline
\textbf{Última revisão} & Joinville, 18 de agosto de 2026 \\
\hline
\textbf{Entrega desta documentação} & 27 de agosto de 2026, até 23h59 \\
\hline
\end{tabular}
\vspace{4pt}

\tableofcontents

% ============================================================
% 1. STATUS ATUAL DO PROJETO
% ============================================================
\section{Status Atual do Projeto}

Este capítulo descreve o estado consolidado do projeto Torque Gestão ao final da migração do material herdado do PAC V e do onboarding inicial da equipe no PAC VI, servindo de linha de base (\emph{baseline}) para o planejamento do próximo MVP.

\subsection{Entregas concluídas}

\begin{itemize}
    \item \textbf{Protótipo de alta fidelidade:} 19 telas mapeadas e navegáveis, desenvolvidas em React (Babel Standalone), publicadas em \url{https://torque-gestao.surge.sh}. Cobrem os três perfis de usuário do sistema --- Administrador, Mecânico e Cliente --- incluindo login unificado, dashboards, gestão completa de ordens de serviço (OS), cadastro de clientes e veículos, catálogo de serviços/peças, gestão de usuários, configurações e portal do cliente.
    \item \textbf{Documentação de especificação (herdada do PAC V):} relatório final com Introdução, Público-alvo, Objetivos, Definição das 19 interfaces, Mapeamento de Riscos (incluindo matriz de riscos), 6 Requisitos Funcionais (RF01--RF06), 6 Requisitos Não Funcionais (RNF01--RNF06) e Stack Tecnológica --- disponível em \texttt{docs/academic/report.pdf}.
    \item \textbf{Diagramas técnicos:} C4 (Contexto, Containers, Componentes), fluxo de processo de abertura de OS, diagrama de sequência, diagrama de estados da OS, modelo Entidade-Relacionamento e diagrama de casos de uso --- todos versionados em \texttt{docs/diagramas/} (Mermaid e PlantUML, com exportações PNG).
    \item \textbf{Design system:} biblioteca de componentes de UI documentada em \texttt{docs/design-system/} (botões, cards, badges, cores, tipografia, inputs, navegação, tabelas, tokens de espaçamento) e prévia consolidada em \texttt{packages/ui/components}.
    \item \textbf{Repositório do PAC VI estruturado:} novo repositório criado (\texttt{torque-gestao}), com monorepo organizado em \texttt{apps/}, \texttt{packages/}, \texttt{docs/}, \texttt{assets/} e \texttt{scripts/}, e pipeline de automação (\texttt{scripts/compile\_report.py}, \texttt{generate\_diagrams.py}, \texttt{generate\_docs.py}) para manter documentação e diagramas sincronizados com o protótipo.
    \item \textbf{Equipe completa como colaboradora do repositório:} os quatro integrantes já têm acesso de escrita ao repositório do PAC VI, superando a pendência registrada no diagnóstico inicial da disciplina.
\end{itemize}

\subsection{O que ainda não foi iniciado}

\begin{itemize}
    \item \textbf{Back-end:} nenhuma linha de código de servidor foi escrita. Não existe ainda o pacote \texttt{apps/api} nem qualquer endpoint FastAPI.
    \item \textbf{Banco de dados:} o PostgreSQL definido na stack ainda não foi provisionado; os dados do protótipo são inteiramente mockados em \texttt{apps/prototype/src/mock-data.jsx}, sem persistência real.
    \item \textbf{Autenticação real:} o login do protótipo simula o acesso por perfil via \texttt{localStorage}, sem validação de credenciais, emissão de token ou controle de acesso efetivo (RF06 ainda não implementado).
    \item \textbf{Integração contínua e deploy:} Docker, GitHub Actions, e as hospedagens em Render e Vercel previstas na stack ainda não foram configuradas.
\end{itemize}

\subsection{Síntese}

Em resumo, o MVP atual do projeto é um \textbf{protótipo de interface completo e navegável, com toda a especificação funcional, os riscos e a arquitetura já documentados}, mas sem nenhuma camada de back-end, persistência ou autenticação real implementada. O próximo MVP, detalhado na seção seguinte, tem como missão transformar esse protótipo estático em uma aplicação funcional de ponta a ponta.

% ============================================================
% 2. PRÓXIMO MVP
% ============================================================
\section{Próximo MVP (MVP 2)}

\subsection{Objetivo}

Evoluir o Torque Gestão de um protótipo de interface com dados mockados para uma \textbf{aplicação funcional de ponta a ponta}, com back-end real em FastAPI, persistência em PostgreSQL e autenticação efetiva, cobrindo o fluxo essencial de uma oficina: cadastro de clientes e veículos, abertura e acompanhamento de ordens de serviço, consulta ao catálogo de serviços e peças, e leitura do status real pelo próprio cliente.

\subsection{Escopo funcional priorizado}

O escopo do MVP 2 foi priorizado a partir dos Requisitos Funcionais já especificados no relatório do PAC V. Os requisitos ligados ao fluxo central da oficina (RF01 a RF06) compõem o núcleo do MVP; funcionalidades administrativas de apoio (relatórios financeiros e KPIs avançados, gestão completa de usuários pela interface) ficam para uma terceira iteração, priorizando primeiro um fluxo de ponta a ponta funcionando de verdade.

\vspace{6pt}
\renewcommand{\arraystretch}{1.4}
\begin{longtable}{|L{1.3cm}|L{4cm}|L{9cm}|}
\caption{Escopo funcional do MVP 2} \\
\hline
\textbf{ID} & \textbf{Requisito} & \textbf{Escopo no MVP 2} \\
\hline
\endfirsthead
\multicolumn{3}{r}{\small(continuação)} \\[4pt]
\hline
\textbf{ID} & \textbf{Requisito} & \textbf{Escopo no MVP 2} \\
\hline
\endhead
\hline
\multicolumn{3}{r}{\small(continua na próxima página)} \\
\endfoot
\hline
\endlastfoot
RF06 & Autenticação e Controle de Acesso & Endpoint de login real, emissão de JWT e middleware de RBAC (Admin, Mecânico, Cliente) protegendo as demais rotas. Implementado primeiro por ser pré-requisito de todo o restante. \\
\hline
RF01 & Gerenciamento de Clientes & CRUD completo de clientes via API, com validação de CPF/CNPJ e busca, substituindo o mock em pelo menos duas telas do protótipo. \\
\hline
RF02 & Cadastro de Veículos & CRUD de veículos vinculados a um cliente, com validação de placa (padrão Mercosul e antigo). \\
\hline
RF03 & Emissão de Ordens de Serviço & Criação de OS com itens de mão de obra e peças separados, cálculo automático do orçamento a partir do catálogo. \\
\hline
RF04 & Acompanhamento de Status & Persistência do fluxo de estados da OS (Aguardando Diagnóstico $\rightarrow$ Em Execução $\rightarrow$ Aguardando Peças $\rightarrow$ Finalizada $\rightarrow$ Entregue), com registro em \texttt{HISTORICO\_STATUS} a cada transição. \\
\hline
RF05 & Catálogo de Serviços e Peças & Endpoints de consulta (e cadastro por Admin) do catálogo, usados na composição do orçamento da OS. \\
\hline
\end{longtable}

\noindent\textbf{Portal do cliente:} passa a ler dados reais no MVP 2. Reaproveitando os mesmos endpoints de RF01, RF02 e RF04 (sem back-end adicional), o cliente autenticado via RF06 consulta o status real de suas OS e o histórico real de seus veículos --- eliminando a lacuna que existia no rascunho anterior deste documento, em que o portal ficava ambíguo entre incluído e não incluído no MVP 2.

\noindent\textbf{Fora do escopo do MVP 2} (reservado para o MVP 3): relatório financeiro e KPIs avançados do dashboard (dados agregados reais), tela de configurações do sistema, e refinamentos de UX do portal do cliente que dependem de infraestrutura adicional --- notificações em tempo real (e-mail/SMS) e atualização automática via \emph{WebSocket/polling}, que exigem componentes fora do escopo de back-end definido nesta iteração.

\subsection{Arquitetura e modelo de dados}

\begin{itemize}
    \item \textbf{Novo pacote \texttt{apps/api}:} serviço FastAPI organizado em camadas de \emph{routers} (endpoints REST), \emph{schemas} (validação Pydantic) e \emph{models} (ORM SQLAlchemy), seguindo a stack já definida no relatório do PAC V.
    \item \textbf{Banco de dados:} PostgreSQL, com schema derivado diretamente do modelo Entidade-Relacionamento já modelado em \texttt{docs/diagramas/modelo-er.md} --- entidades \texttt{CLIENTE}, \texttt{VEICULO}, \texttt{USUARIO}, \texttt{ORDEM\_SERVICO}, \texttt{ITEM\_OS}, \texttt{HISTORICO\_STATUS}, \texttt{CATALOGO\_SERVICO} e \texttt{CATALOGO\_PECA} --- versionado com migrations (Alembic).
    \item \textbf{Autenticação:} tokens JWT com claim de perfil (\texttt{role}), validados por \emph{dependency injection} do FastAPI em cada rota protegida, implementando o RBAC previsto em RF06/RNF01.
    \item \textbf{Documentação de API:} geração automática de especificação OpenAPI/Swagger (RNF06), publicada em \texttt{/docs} do serviço, servindo como coleção de referência para o time de front-end e como contrato formal entre back-end e front-end (ver mitigação do risco de inconsistência de contrato, Seção~\ref{sec:riscos-mvp2}).
    \item \textbf{Front-end:} as telas prioritárias do protótipo (login, dashboard admin, lista/nova OS, detalhes da OS, cadastro de clientes e veículos, e as três telas do portal do cliente) passam a consumir a API real via \texttt{fetch}, substituindo progressivamente \texttt{apps/prototype/src/mock-data.jsx}.
\end{itemize}

\subsection{Entregáveis técnicos}

\begin{itemize}
    \item Pacote \texttt{apps/api} com os endpoints de RF01, RF02, RF03, RF04, RF05 e RF06 implementados e testados.
    \item Script de migrations (Alembic) e \emph{seed} inicial do catálogo de serviços/peças.
    \item \texttt{Dockerfile} do back-end e \texttt{docker-compose.yml} (API + PostgreSQL) para ambiente de desenvolvimento local, conforme a stack de conteinerização já definida.
    \item Workflow de GitHub Actions executando lint e suíte de testes automatizados (\texttt{pytest}) a cada \emph{pull request}.
    \item Suíte de testes cobrindo as regras de negócio críticas: cálculo de orçamento (RF03) e transições de status da OS (RF04).
    \item Telas do protótipo (login, gestão de OS, cadastro de clientes/veículos e portal do cliente) conectadas à API real.
    \item Atualização do README com instruções de execução local (\emph{passo a passo}) e, caso o schema real divirja do modelo mockado, atualização de \texttt{docs/diagramas/modelo-er.md}.
\end{itemize}

\subsection{Critérios de aceite (Definition of Done do MVP 2)}

O MVP 2 será considerado concluído quando, cumulativamente:

\begin{enumerate}
    \item o ambiente completo (API + banco de dados) subir localmente com um único comando (\texttt{docker compose up});
    \item o login autenticar de fato contra o banco de dados, emitindo um JWT válido e respeitando o RBAC por perfil;
    \item cadastro de clientes e veículos persistir no PostgreSQL, sobrevivendo a um \emph{reload} da aplicação;
    \item a criação e a atualização de status de uma OS persistirem e respeitarem o fluxo de estados definido em RF04;
    \item o catálogo de serviços e peças for consultável via API e utilizado no cálculo do orçamento da OS;
    \item a documentação OpenAPI/Swagger estiver publicada e acessível;
    \item as telas de login, gestão de OS, cadastro de clientes/veículos \textbf{e o portal do cliente (leitura de status e histórico)} consumirem dados reais da API, em vez do mock;
    \item o pipeline de CI (lint + testes) estiver executando automaticamente a cada \emph{pull request}.
\end{enumerate}

\subsection{Cronograma e divisão de responsabilidades}
\label{sec:cronograma-mvp2}

O professor orientador do PAC VI definiu que, até \textbf{27 de agosto de 2026}, além desta documentação, a equipe deve ter \textbf{uma parte relevante do MVP 2 efetivamente implementada} --- não apenas planejada. Como restam poucos dias úteis a partir de hoje (\textbf{18 de agosto}), as antigas Fase 1 (infraestrutura) e Fase 2 (identidade e cadastros base) foram compactadas em uma \textbf{Sprint 1}, com os quatro integrantes trabalhando em paralelo desde já, em vez de em sequência. Isso só é possível porque o contrato de dados já está fechado no modelo Entidade-Relacionamento (\texttt{docs/diagramas/modelo-er.md}): cada integrante pode começar sua parte imediatamente, sem esperar a infraestrutura do Miguel Angelo estar pronta, e integrar ao PostgreSQL real assim que ela existir (por volta do dia 21/08).

\vspace{6pt}
\begingroup
\small
\renewcommand{\arraystretch}{1.4}
\begin{longtable}{|L{4.1cm}|L{10.9cm}|}
\caption{Sprint 1 --- cronograma diário por integrante (18 a 27/08)} \label{tab:sprint1} \\
\hline
\textbf{Responsável} & \textbf{Etapas e datas} \\
\hline
\endfirsthead
\multicolumn{2}{r}{\small(continuação)} \\[4pt]
\hline
\textbf{Responsável} & \textbf{Etapas e datas} \\
\hline
\endhead
\hline
\multicolumn{2}{r}{\small(continua na próxima página)} \\
\endfoot
\hline
\endlastfoot
Miguel Angelo Dufloth Filho \newline (infraestrutura) &
\begin{itemize}[leftmargin=1em, itemsep=2pt, topsep=2pt]
    \item \textbf{18--19/08:} pacote \texttt{apps/api} (camadas \emph{routers}/\emph{schemas}/\emph{models}) e \texttt{docker-compose.yml} (API + PostgreSQL).
    \item \textbf{20--21/08:} \emph{models} SQLAlchemy de todas as entidades do modelo ER; Alembic configurado e primeira \emph{migration} aplicada.
    \item \textbf{22--24/08:} apoio à integração dos demais integrantes contra o PostgreSQL real.
    \item \textbf{25--27/08:} validação de ponta a ponta (\texttt{docker compose up} com todos os módulos plugados) e ajustes finais.
\end{itemize} \\
\hline
Miguel Angel Balladares Huertas \newline (RF06 --- autenticação) &
\begin{itemize}[leftmargin=1em, itemsep=2pt, topsep=2pt]
    \item \textbf{18--20/08:} endpoint de login, \emph{hash} de senha com Bcrypt e emissão de JWT (testado localmente, sem depender do Postgres).
    \item \textbf{21--23/08:} \emph{middleware}/\emph{dependency} de RBAC (Admin, Mecânico, Cliente) protegendo rotas.
    \item \textbf{24--26/08:} integração com o PostgreSQL real assim que a infraestrutura estiver pronta; testes manuais via Swagger.
    \item \textbf{27/08:} ajustes finais e demonstração.
\end{itemize} \\
\hline
Leonardo Lotério de Lima \newline (RF01 --- clientes) &
\begin{itemize}[leftmargin=1em, itemsep=2pt, topsep=2pt]
    \item \textbf{18--20/08:} \emph{schema} Pydantic e \emph{model} de Cliente; CRUD básico em ambiente local.
    \item \textbf{21--23/08:} validação de CPF/CNPJ e busca por nome ou documento.
    \item \textbf{24--26/08:} integração com o PostgreSQL real; testes manuais via Swagger.
    \item \textbf{27/08:} ajustes finais e demonstração.
\end{itemize} \\
\hline
Lucas Honorato dos Santos \newline (RF02 --- veículos) &
\begin{itemize}[leftmargin=1em, itemsep=2pt, topsep=2pt]
    \item \textbf{18--20/08:} \emph{schema} Pydantic e \emph{model} de Veículo, acompanhando a estrutura de Cliente do Leonardo (relação de chave estrangeira).
    \item \textbf{21--23/08:} CRUD de veículos vinculado a um cliente; validação de placa (padrão Mercosul e antigo).
    \item \textbf{24--26/08:} integração com o PostgreSQL real; testes manuais via Swagger.
    \item \textbf{27/08:} ajustes finais, demonstração e checklist do que foi entregue para reportar ao professor.
\end{itemize} \\
\hline
\end{longtable}
\endgroup

\noindent\textbf{Entrega da Sprint 1 (27/08):} com um único comando (\texttt{docker compose up}), a API e o PostgreSQL sobem localmente; o login autentica de fato e emite um JWT válido; o RBAC protege as rotas por perfil; e o CRUD de clientes e de veículos persiste no banco de dados real --- tudo demonstrável via Swagger. RF03/RF04 (ordens de serviço) e RF05 (catálogo) dependem de clientes e veículos já existirem, por isso permanecem para a fase seguinte, e não porque ficaram de fora por atraso.

\vspace{6pt}
\begingroup
\small
\renewcommand{\arraystretch}{1.5}
\begin{longtable}{|L{4.1cm}|L{3cm}|L{2cm}|L{4.9cm}|}
\caption{Fases seguintes do MVP 2 (a partir de 28/08)} \label{tab:fases-seguintes} \\
\hline
\textbf{Fase} & \textbf{Responsável} & \textbf{Período} & \textbf{Entrega da fase} \\
\hline
\endfirsthead
\multicolumn{4}{r}{\small(continuação)} \\[4pt]
\hline
\textbf{Fase} & \textbf{Responsável} & \textbf{Período} & \textbf{Entrega da fase} \\
\hline
\endhead
\hline
\multicolumn{4}{r}{\small(continua na próxima página)} \\
\endfoot
\hline
\endlastfoot
Fase 3 --- Núcleo operacional & Leonardo Lotério de Lima & 28/08 a 10/09 & RF05 (catálogo) e RF03/RF04 (emissão de OS e histórico de status) implementados, com cálculo de orçamento a partir do catálogo. \\
\hline
Fase 4 --- Integração front-end e portal do cliente & Lucas Honorato dos Santos & 11/09 a 24/09 & Telas de login, gestão de OS, cadastro de clientes/veículos e portal do cliente consumindo a API real; suíte de testes automatizados das regras críticas de RF03/RF04. \\
\hline
Fase 5 --- CI/CD e deploy de validação & Miguel Angelo Dufloth Filho e Lucas Honorato dos Santos & 25/09 a 01/10 & Pipeline de GitHub Actions (lint + \texttt{pytest}) ativo a cada \emph{pull request}; deploy de validação em Render (API/BD) e Vercel (front-end). \\
\hline
\end{longtable}
\endgroup

\noindent\textit{Observação: as datas da Sprint 1 são o compromisso firme da equipe para a entrega de 27/08. As datas das fases seguintes são o planejamento atual e ficam sujeitas a ajuste caso o professor orientador do PAC VI sinalize um calendário oficial diferente para os marcos do MVP 2.}

\subsection{Acompanhamento e comunicação}

Para reduzir o risco de atrasos silenciosos --- já identificado na Seção~\ref{sec:riscos-mvp2} como sobreposição de calendário com outras disciplinas --- a equipe adota:

\begin{itemize}
    \item \textbf{Reunião semanal de status}, toda sexta-feira, com atualização do progresso de cada fase e realocação de tarefas entre os quatro integrantes caso alguma fase esteja atrasada. Durante a Sprint 1 (18 a 27/08), essa checagem é diária, dado o prazo curto e o compromisso assumido com o professor.
    \item \textbf{Quadro de acompanhamento no GitHub} (Issues/Projects), com uma \emph{issue} por requisito (RF01--RF06) e por entregável técnico da Seção~2.4, atribuída ao responsável correspondente nas Tabelas~\ref{tab:sprint1} e~\ref{tab:fases-seguintes}.
    \item Cada responsável atualiza o status da própria \emph{issue} até a véspera da reunião de acompanhamento, permitindo identificar bloqueios antes que afetem a fase seguinte.
\end{itemize}

\subsection{Riscos específicos do MVP 2}
\label{sec:riscos-mvp2}

\vspace{6pt}
\begingroup
\small
\renewcommand{\arraystretch}{1.5}
\begin{longtable}{|L{3cm}|L{2.1cm}|L{0.9cm}|L{1cm}|L{0.9cm}|L{5cm}|}
\hline
\textbf{Risco} & \textbf{Categoria} & \textbf{Prob.} & \textbf{Impacto} & \textbf{Nível} & \textbf{Mitigação} \\
\hline
\endfirsthead
\multicolumn{6}{r}{\small(continuação)} \\[4pt]
\hline
\textbf{Risco} & \textbf{Categoria} & \textbf{Prob.} & \textbf{Impacto} & \textbf{Nível} & \textbf{Mitigação} \\
\hline
\endhead
\hline
\multicolumn{6}{r}{\small(continua na próxima página)} \\
\endfoot
\hline
\endlastfoot
Inconsistências de contrato entre API real e front-end mockado & Técnico & Média & Médio & Médio & Contrato definido via OpenAPI/Swagger já na Sprint 1; front-end consome o schema gerado automaticamente em vez de assumir formatos manualmente. \\
\hline
Curva de aprendizado da equipe com FastAPI, PostgreSQL e Alembic & Organizacional & Média & Médio & Médio & Sessão interna de \emph{pair programming} no início da Sprint 1, com tutorial guiado compartilhado entre os quatro integrantes desde o primeiro dia. \\
\hline
Sobreposição de calendário com entregas de outra disciplina do semestre & Organizacional & Média & Baixo & Baixo & Reuniões semanais fixas (Seção 2.7) para identificar atrasos cedo e redistribuir tarefas entre os quatro integrantes quando necessário. \\
\hline
Atraso na configuração de infraestrutura de deploy (Render/Vercel) & Técnico & Baixa & Médio & Baixo & Fase 5 isolada como etapa final sem dependências pendentes de outras fases; o deploy de validação pode escorregar sem travar as entregas anteriores do MVP 2. \\
\hline
Riscos de segurança já mapeados no PAC V (SQL Injection, falhas de autenticação) migrarem do papel para a implementação sem a devida mitigação & Segurança & Baixa & Alto & Médio & Checklist de segurança (Bcrypt, JWT, ORM, HTTPS) revisado como critério de aceite em cada \emph{pull request}, não apenas ao final do MVP 2. \\
\hline
\end{longtable}
\endgroup

\end{document}
